\lecture{9.5}{Instrumental Variables}
\label{lec:9-5}

\subsection{IV Graph}

An instrumental variable is a variable $Z$ that affects treatment $A$ but is
otherwise unrelated to the outcome $Y$ except through $A$. It is useful when
there is unmeasured confounding between $A$ and $Y$.

\begin{figure}[H]
\centering
\begin{tikzpicture}[
  node distance=1.35cm,
  box/.style={draw, rounded corners=2pt, minimum width=.9cm, minimum height=.7cm},
  arr/.style={-{Latex[length=2.3mm]}, thick},
  bi/.style={<->, >=Latex, thick, dashed}
]
\node[box] (Z) {$Z$};
\node[box, right=of Z] (A) {$A$};
\node[box, right=of A] (Y) {$Y$};
\node[box, above=of A] (U) {$U$};
\draw[arr] (Z) -- (A);
\draw[arr] (A) -- (Y);
\draw[arr] (U) -- (A);
\draw[arr] (U) -- (Y);
\end{tikzpicture}
\caption{Instrumental variable graph: $Z$ affects $Y$ only through $A$, while $U$ confounds $A$ and $Y$.}
\label{fig:ch9-iv-graph}
\end{figure}

\subsection{Relevance}

\begin{assumptionbox}
The instrument must be associated with treatment:
\[
  Z \nci A.
\]
In linear models this is often expressed as
\[
  \Cov(Z,A)\ne 0.
\]
\end{assumptionbox}

Relevance is partly empirically checkable because $Z$ and $A$ are observed.
Weak relevance leads to unstable IV estimates.

\subsection{Marginal Ignorability}

\begin{assumptionbox}
The instrument is marginally independent of the potential outcomes:
\[
  Y(a,z)\ci Z
  \qquad\text{for all }a,z.
\]
\end{assumptionbox}

This says $Z$ is as-if randomized with respect to the outcome process. In
practice, this is a scientific assumption: an observed association between $Z$
and baseline covariates may cast doubt on it, but no observed-data test can
prove it.

\subsection{Exclusion Restriction}

\begin{assumptionbox}
The instrument has no direct effect on the outcome:
\[
  Y(a,z)=Y(a)
  \qquad\text{for all }a,z.
\]
\end{assumptionbox}

The exclusion restriction rules out all pathways from $Z$ to $Y$ except through
$A$. It is often the most controversial IV assumption.

\subsection{Linear Model Assumption}

In the simplest linear structural model, suppose variables are standardized and
\[
  A=\gamma_{az}Z+\gamma_{au}U+\epsilon_A,
\]
\[
  Y=\gamma_{ya}A+\gamma_{yu}U+\epsilon_Y.
\]
The causal effect of $A$ on $Y$ is $\gamma_{ya}$. Because $U$ is unobserved,
the ordinary association between $A$ and $Y$ is confounded.

\subsection{IV Estimand}

Under the IV assumptions and the linear model,
\[
  \Cov(Z,Y)=\gamma_{ya}\Cov(Z,A).
\]
Therefore
\[
  \gamma_{ya}=\frac{\Cov(Z,Y)}{\Cov(Z,A)}.
\]
If variables are standardized, covariances can be written as correlations:
\[
  \gamma_{ya}=\frac{\operatorname{Corr}(Z,Y)}
  {\operatorname{Corr}(Z,A)}.
\]

\subsection{Wald Ratio Intuition}

For a binary instrument, the Wald estimand is
\[
  \frac{\E(Y\given Z=1)-\E(Y\given Z=0)}
       {\E(A\given Z=1)-\E(A\given Z=0)}.
\]
The numerator is the effect of the instrument on the outcome. The denominator is
the effect of the instrument on the treatment. If the instrument affects the
outcome only by shifting treatment, the ratio rescales the intention-to-treat
effect by the amount of treatment movement induced by the instrument.
With a binary instrument and a binary or otherwise nonlinear treatment setting,
the Wald ratio is commonly interpreted as a LATE or CACE only under additional
assumptions such as monotonicity, exclusion, and appropriate independence
assumptions; without them, it is not automatically the population ATE.

\begin{identificationbox}
In a linear IV model with relevance, instrument ignorability, and exclusion,
\[
  \text{causal effect of }A\text{ on }Y
  =
  \frac{\Cov(Z,Y)}{\Cov(Z,A)}.
\]
\end{identificationbox}

\subsection{Limitations of IV Analysis}

IV analysis is powerful but fragile. Limitations include:
\begin{itemize}
  \item exclusion restrictions are untestable and often scientifically strong;
  \item weak instruments produce unstable estimates;
  \item linear IV identifies a model-specific effect, not necessarily a
  nonparametric ATE;
  \item with heterogeneous treatment effects, additional assumptions are needed
  to interpret the estimand as an effect for a particular subpopulation;
  \item invalid instruments can amplify rather than reduce bias.
\end{itemize}

\begin{keyidea}
An instrument replaces adjustment for unmeasured confounding with different
untestable assumptions: the instrument must move treatment, be as-if randomized,
and affect the outcome only through treatment.
\end{keyidea}
